\section{Casi d'Uso}

I seguenti casi d'uso descrivono in modo concreto e passo-passo le interazioni che gli utenti di Culturia intraprenderanno con il sistema. Ogni scenario è ancorato ai sotto-task definiti nell'Assignment~1, garantendo coerenza con l'analisi dei bisogni emersi dalla ricerca sugli utenti.

% ─── UC1 ─────────────────────────────────────────────────────────
\subsection*{UC1 — Scansione RFID e accesso ai contenuti del luogo}

\noindent
\begingroup
\renewcommand{\tabularxcolumn}[1]{m{#1}}
\renewcommand{\arraystretch}{1.6}
\begin{tabularx}{\textwidth}{@{} >{\bfseries\color{culturiaRed}}m{2.2cm} X @{}}
  \toprule
  Persona     & Marco, 46 anni --- Turista Esterno \\
  Sotto-task  & 1.1 Scansione RFID \\
  Contesto    & Marco arriva per la prima volta in un polo culturale. All'ingresso nota un pannello informativo con il logo di Culturia e le istruzioni per avvicinare il proprio smartphone. Non ha mai usato l'applicazione prima d'ora. \\
  \bottomrule
\end{tabularx}
\endgroup

\vspace{0.3cm}
\noindent\textbf{\color{culturiaRed}Sequenza di azioni:}
\begin{enumerate}[nosep, leftmargin=1.4em]
  \item Marco avvicina lo smartphone al tag RFID integrato nel pannello informativo all'ingresso.
  \item L'applicazione Culturia si attiva automaticamente, rileva il luogo culturale e ne carica i contenuti contestuali (nome, descrizione breve, orari e mappa interna).
  \item Sulla schermata principale compare un pannello di benvenuto con i contenuti del luogo e un menu di azioni disponibili:
    \begin{itemize}[nosep, leftmargin=1.6em]
      \item \textbf{Tour in Realtà Aumentata} — esplora le opere con overlay informativi;
      \item \textbf{Spiegazione AI} — ascolta la descrizione dell'opera tramite il Cicerone intelligente;
      \item \textbf{Descrizione testuale} — leggi le informazioni storiche e artistiche del luogo;
      \item \textbf{Strumenti di accessibilità} — attiva percorsi senza barriere, sottotitoli o testo ingrandito.
    \end{itemize}
  \item Marco scorre il pannello, valuta le opzioni disponibili e sceglie da quale funzionalità iniziare la sua visita.
\end{enumerate}

\vspace{0.15cm}
\noindent\textbf{\color{culturiaRed}Risultato atteso:}
Marco ottiene in pochi secondi una panoramica completa del luogo e di tutte le funzionalità che l'app mette a sua disposizione, senza dover navigare manualmente tra menu complessi. L'esperienza di ingresso è immediata, intuitiva e orientata all'azione.

\vspace{0.5cm}

% ─── UC2 ─────────────────────────────────────────────────────────
\subsection*{UC2 — Scansione visiva ed overlay in Realtà Aumentata}

\noindent
\begingroup
\renewcommand{\tabularxcolumn}[1]{m{#1}}
\renewcommand{\arraystretch}{1.6}
\begin{tabularx}{\textwidth}{@{} >{\bfseries\color{culturiaRed}}m{2.2cm} X @{}}
  \toprule
  Persona     & Elena, 26 anni --- Turista Interna (Visitatrice Assidua) \\
  Sotto-task  & 2.1 Scansione visiva ad overlay (AR), 2.2 Selezione Tone of Voice (IA), 2.3 Interrogazione vocale dinamica \\
  Contesto    & Elena ha appena completato la scansione RFID (UC1) ed è all'interno del polo culturale. Decide di avviare il tour in Realtà Aumentata per approfondire uno degli affreschi presenti nella sala nord. \\
  \bottomrule
\end{tabularx}
\endgroup

\vspace{0.3cm}
\noindent\textbf{\color{culturiaRed}Sequenza di azioni:}
\begin{enumerate}[nosep, leftmargin=1.4em]
  \item Partendo dal pannello di benvenuto (UC1), Elena tocca il pulsante \textbf{Tour in Realtà Aumentata}.
  \item L'app apre la fotocamera in modalità AR e invita l'utente a inquadrare le opere circostanti.
  \item Elena punta la fotocamera verso l'affresco medievale: il sistema lo riconosce automaticamente e sovrappone un overlay con le informazioni contestuali più rilevanti:
    \begin{itemize}[nosep, leftmargin=1.6em]
      \item nome e datazione dell'opera;
      \item annotazioni evidenziate sulle aree di restauro;
      \item icone interattive per approfondire specifici dettagli.
    \end{itemize}
  \item Nell'interfaccia AR è visibile anche la \textbf{mappa del percorso consigliato}: frecce direzionali indicano le tappe successive del tour, aiutando Elena a navigare nella sala senza perdere il filo del percorso.
  \item Elena seleziona il Tone of Voice ``Accademico'' per ricevere descrizioni con terminologia specialistica, adatta al suo percorso universitario.
  \item Incuriosita da una figura secondaria nell'angolo inferiore, formula una domanda vocale al Cicerone IA per ottenere informazioni più verticali sul soggetto rappresentato.
  \item Il sistema risponde con una spiegazione dettagliata e suggerisce un'opera correlata in un'altra sala, invitando Elena a proseguire il percorso.
\end{enumerate}

\vspace{0.15cm}
\noindent\textbf{\color{culturiaRed}Risultato atteso:}
Elena fruisce di uno strato informativo ricco e interattivo direttamente sull'opera, senza mai interrompere il flusso della visita. La combinazione di overlay AR, navigazione integrata e dialogo con il Cicerone trasforma la sua esperienza in uno studio attivo dell'opera.

\vspace{0.5cm}

% ─── UC3 ─────────────────────────────────────────────────────────
\subsection*{UC3 — Liste di missioni contestuali}

\noindent
\begingroup
\renewcommand{\tabularxcolumn}[1]{m{#1}}
\renewcommand{\arraystretch}{1.6}
\begin{tabularx}{\textwidth}{@{} >{\bfseries\color{culturiaRed}}m{2.2cm} X @{}}
  \toprule
  Persona     & Matteo, 34 anni --- Turista Interno (Prima Visita) \\
  Sotto-task  & 3.1 Dashboard e progressioni, 3.2 Rewarding tramite tap RFID, 3.3 Liste di missioni contestuali \\
  Contesto    & Matteo ha selezionato il polo culturale dall'elenco dei luoghi in-app ed è entrato nella sezione dedicata. Curioso ma timoroso di annioarsi, decide di esplorare le attività gamificate proposte dall'applicazione. \\
  \bottomrule
\end{tabularx}
\endgroup

\vspace{0.3cm}
\noindent\textbf{\color{culturiaRed}Sequenza di azioni:}
\begin{enumerate}[nosep, leftmargin=1.4em]
  \item Matteo, dopo aver selezionato il polo culturale dalla schermata principale dell'app, accede alla sezione \textbf{Missioni}.
  \item Visualizza una lista di missioni contestuali, suddivise in categorie:
    \begin{itemize}[nosep, leftmargin=1.6em]
      \item \textbf{Tour} — sequenze di opere da visitare in un determinato ordine tematico;
      \item \textbf{Tag da scansionare} — punti RFID nascosti o sfidanti da raggiungere nel sito;
      \item \textbf{Azioni in-app} — obiettivi relativi all'uso delle funzionalità (es.\ ascoltare 3 spiegazioni AI, usare il filtro accessibilità, ecc.).
    \end{itemize}
  \item Matteo seleziona la missione \textit{``Il Drago Nascosto''} e consulta le tappe richieste e il premio digitale previsto al completamento.
  \item Guidato dall'app, raggiunge il primo tag RFID della missione: avvicina lo smartphone e sblocca il primo step, ricevendo un feedback visivo animato e un badge parziale nella dashboard.
  \item Completa le tappe successive, vedendo la propria barra di progressione avanzare in tempo reale.
  \item Al termine, la missione viene segnata come completata: vengono assegnati il badge definitivo, i punti e viene visualizzata la lista delle missioni ancora aperte per incentivare il ritorno.
\end{enumerate}

\vspace{0.15cm}
\noindent\textbf{\color{culturiaRed}Risultato atteso:}
Matteo supera la soglia dell'apatia iniziale grazie alla struttura narrativa e al meccanismo di ricompensa progressiva. La visita diventa un percorso a tappe motivante che lo spinge a esplorare zone del polo che non avrebbe mai visitato autonomamente, e a tornare per completare le missioni rimaste aperte.

\vspace{0.5cm}

% ─── UC4 ─────────────────────────────────────────────────────────
\subsection*{UC4 — Configurazione di un percorso AR guidato}

\noindent
\begingroup
\renewcommand{\tabularxcolumn}[1]{m{#1}}
\renewcommand{\arraystretch}{1.6}
\begin{tabularx}{\textwidth}{@{} >{\bfseries\color{culturiaRed}}m{2.2cm} X @{}}
  \toprule
  Persona     & Giovanni, 55 anni --- Impiegato del polo culturale \\
  Sotto-task  & 4.2 Configurazione percorso AR guidato \\
  Contesto    & Giovanni deve configurare un nuovo tour guidato in AR per una mostra temporanea appena allestita. Le schede delle opere sono già presenti nella piattaforma. Accede all'interfaccia di authoring di Culturia, progettata per essere usata in autonomia da figure non tecniche, in modo simile a strumenti come NotebookLM. \\
  \bottomrule
\end{tabularx}
\endgroup

\vspace{0.3cm}
\noindent\textbf{\color{culturiaRed}Sequenza di azioni:}
\begin{enumerate}[nosep, leftmargin=1.4em]
  \item Giovanni accede al pannello di amministrazione e apre la sezione \textbf{Percorsi AR}, quindi seleziona \textbf{Crea nuovo percorso}.
  \item Nella prima schermata, carica le fonti di riferimento per il tour: documenti PDF, cataloghi della mostra e immagini delle opere, che alimentano la base di conoscenza del Cicerone IA.
  \item Inserisce un \textbf{prompt testuale} per istruire l'IA su come narrare il percorso (es.\ \textit{``Tono avventuroso per un pubblico giovane, con focus sul simbolismo dei colori''}).
  \item Accede alla \textbf{mappa interattiva} del polo culturale e, tramite drag-and-drop, seleziona e ordina i punti di interesse che comporranno le tappe del tour.
  \item Per ciascuna tappa, associa il tag RFID fisico corrispondente all'opera o all'area espositiva tramite un menù a tendina.
  \item Avvia la \textbf{modalità anteprima}: visualizza il percorso simulato con le frecce AR e i testi generati dall'IA, verificando la coerenza narrativa e spaziale del tour.
  \item Apporta eventuali aggiustamenti (riordina le tappe, modifica il prompt) e preme \textbf{Pubblica}: il tour diventa immediatamente disponibile per i visitatori nell'app.
\end{enumerate}

\vspace{0.15cm}
\noindent\textbf{\color{culturiaRed}Risultato atteso:}
Giovanni configura e pubblica un tour AR completo in totale autonomia, senza dover coinvolgere tecnici esterni. L'interfaccia visuale con mappa, prompt testuale e anteprima integrata abbatte la barriera tecnica, consentendo a figure non specializzate di produrre percorsi culturali di qualità professionale.
